\subsection{Plan comercial con estimación de ventas}

\subsubsection{Definición del producto}

El middleware propuesto es una solución que actúa como intermediario entre las fuentes de datos de una empresa y sus asistentes virtuales, proporcionando contexto en tiempo real para generar respuestas hiperpersonalizadas. Se vende en forma de licencias con renovación anual. La implementación se cobra aparte y es realizada por un equipo de consultores propios.

\subsubsection{Segmentación del mercado objetivo}

Se identifican los siguientes sectores como los más beneficiados por la implementación del middleware:

\begin{itemize}
    \item Banca y Servicios Financieros
    \item Comercio Minorista y Comercio Electrónico
    \item Atención Médica y Ciencias de la Vida
    \item Telecomunicaciones y Tecnología de la Información
    \item Viajes y Hospitalidad
\end{itemize}

\subsubsection{Estrategia de penetración en el mercado}

Para cada sector, se desarrollarán estrategias específicas que incluyen:

\begin{itemize}
    \item Desarrollar soluciones que aborden problemas particulares de cada industria.
    \item Asegurar que el middleware cumpla con las regulaciones específicas de cada sector.
    \item Integración con los sistemas y plataformas existentes en cada industria. Principalmente CRMs.
\end{itemize}

\subsubsection{Modelo de ingresos}

Se propone un modelo de ingresos basado en licencias de software con el cobro de una tarifa inicial por la implementación del middleware y renovaciones anuales.

Para fundamentar una proyección de ventas sólida para el middleware de hiperpersonalización en asistentes virtuales, es esencial analizar el mercado objetivo, sus tendencias de crecimiento y las oportunidades específicas en cada sector.

\subsubsection{Tamaño y crecimiento del mercado de asistentes virtuales}

Según \textcite{emergen_virtual_assistants_2023}, el mercado global de asistentes virtuales ha mostrado un crecimiento significativo en los últimos años y se espera que esta tendencia continúe:

\begin{itemize}
    \item Valor de Mercado en 2022: USD 11,40 mil millones.
    \item Proyección para 2032: USD 150,58 mil millones.
    \item Tasa de Crecimiento Anual Compuesta (CAGR): 29,6\% durante el período 2022-2032.
\end{itemize}

\subsubsection{Mercado de hiperpersonalización}

También \textcite{emergen_hyperpersonalization_2024} afirma que la hiperpersonalización está en expansión:

\begin{itemize}
    \item Valor de Mercado en 2024: USD 19,37 mil millones.
    \item Tasa de Crecimiento Anual Compuesta Proyectada: 15,83\% durante el período de pronóstico.
\end{itemize}

\subsubsection{Adopción de asistentes virtuales en américa latina}

De acuerdo a un reporte de \textcite{informes_expertos_latam_2024}, en América Latina, la adopción de asistentes virtuales inteligentes está creciendo a tasas similares:

\begin{itemize}
    \item Valor de Mercado en 2023: USD 5,43 mil millones.
    \item Proyección para 2032: USD 16,18 mil millones.
    \item Tasa de Crecimiento Anual Compuesta Proyectada: 12,9\% entre 2024 y 2032.
\end{itemize} 